\documentclass[11pt,a4paper]{article}
\usepackage[T1]{fontenc}
\usepackage{isabelle,isabellesym}

% This should be the last package used.
\usepackage{pdfsetup}

\urlstyle{rm}
\isabellestyle{it}

\begin{document}

\title{The Teichm\"uller--Tukey Lemma}
\author{%
  Vithor Lindermann Kraisch\\
  \small Department of Mobility Engineering\\
  \small Federal University of Santa Catarina, Joinville, Brazil
  \and
  Luiz Gustavo Cordeiro\\
  \small Department of Mobility Engineering\\
  \small Federal University of Santa Catarina, Joinville, Brazil}
\maketitle

\begin{abstract}
  This entry formalizes the Teichm\"uller--Tukey lemma in Isabelle/HOL:
  every nonempty family of sets of finite character has a member that is
  maximal under inclusion.  Instead of deriving the result from Zorn's lemma,
  which is already available in Isabelle/HOL, the development follows the
  direct choice-function construction of Sun and Yu, originally formulated in
  Morse--Kelley set theory and checked in Coq.  The Isabelle proof makes
  explicit the choice argument that turns a fixed point of the construction
  into a maximal member.
\end{abstract}

\tableofcontents

\input{session}

\bibliographystyle{abbrv}
\bibliography{root}

\end{document}
